\section{CRDT and Merge Semantics}
\label{sec:CrdtAndMerge}

Archimesh uses a hybrid consistency model: the server is the single point of
serialization for revision ordering, but per-field CRDT merges handle
concurrent edits on different fields of the same object without requiring
explicit locks.

\subsection{Causal identity}

Every operation carries a \texttt{causal} object:

\begin{lstlisting}[caption={Causal metadata}]
{
  "clientId": "<stable client/session id>",
  "lamport": <monotonic per-client counter>,
  "opId": "<idempotency identity>"
}
\end{lstlisting}

If the client omits \texttt{causal}, the server fills it deterministically
from the assigned revision range. The comparison rule for Last-Writer-Wins:

\begin{quote}
\texttt{(lamportA,{} clientIdA)} wins over \texttt{(lamportB,{} clientIdB)} iff
\texttt{lamportA > lamportB}, or \texttt{lamportA == lamportB} and
\texttt{clientIdA >= clientIdB} (lexicographic).
\end{quote}

\subsection{Merge types by domain}

\subsubsection{Semantic element fields}

Fields: \texttt{name}, \texttt{documentation}.

\begin{itemize}
\item Type: LWW Register per field per element.
\item Metadata: field clock map \texttt{elementId -> \{ field -> (lamport, clientId) \}}.
\item Merge: apply incoming field if the incoming tuple wins the field clock;
  otherwise ignore.
\end{itemize}

\subsubsection{Relationship fields}

Fields: \texttt{name}, \texttt{documentation}, \texttt{sourceId},
\texttt{targetId}.

Same LWW register semantics as elements, with per-field clocks stored on the
\texttt{:Relationship} node.

\subsubsection{Diagram notation fields}

View objects and connections use per-field LWW registers for:

\begin{description}
\item[View objects] \texttt{x}, \texttt{y}, \texttt{width}, \texttt{height},
  \texttt{type}, \texttt{alpha}, \texttt{lineAlpha}, \texttt{lineWidth},
  \texttt{lineStyle}, \texttt{textAlignment}, \texttt{textPosition},
  \texttt{gradient}, \texttt{iconVisibleState},
  \texttt{deriveElementLineColor}, \texttt{fillColor}, \texttt{lineColor},
  \texttt{font}, \texttt{fontColor}, \texttt{iconColor},
  \texttt{imagePath}, \texttt{imagePosition}, \texttt{name},
  \texttt{documentation}.
\item[Connections] \texttt{type}, \texttt{nameVisible},
  \texttt{textAlignment}, \texttt{textPosition}, \texttt{lineWidth},
  \texttt{name}, \texttt{lineColor}, \texttt{font}, \texttt{fontColor},
  \texttt{documentation}, \texttt{bendpoints}.
\end{description}

Per-field notation clocks are stored as properties on the
\texttt{:ViewObject} and \texttt{:Connection} nodes (e.g.
\texttt{notation\_lamport\_x}, \texttt{notation\_clientId\_x}).

\subsubsection{View object child membership}

Type: OR-set. Operations: \texttt{AddViewObjectChildMember},
\texttt{RemoveViewObjectChildMember}.

Member clocks are stored via \texttt{:PropertyClock} nodes in the
\texttt{VIEWOBJECT\_CHILD} namespace. Materialized as
\texttt{(:ViewObject)-[:CHILD\_MEMBER]->(:ViewObject)} edges.

\subsubsection{Property set membership}

Type: OR-set. Operations: \texttt{AddPropertySetMember},
\texttt{RemovePropertySetMember}.

Member clocks are stored via \texttt{:PropertyClock} nodes.
Materialized as canonical JSON string arrays in the snapshot/export output.

\subsection{Delete semantics --- tombstones}

When an entity is deleted, a tombstone is recorded with the deleting
operation's causal clock:

\begin{lstlisting}[caption={Tombstone record}]
{ modelId, id, lamport, clientId }
\end{lstlisting}

Rules:

\begin{itemize}
\item \texttt{Delete*} always records or updates the tombstone (even if the
  entity was already missing).
\item Stale \texttt{Update*} or \texttt{Create*} (with a causal clock older
  than the tombstone) is ignored --- preventing resurrection.
\item A newer \texttt{Create*} (with a causal clock that beats the tombstone)
  may clear the tombstone and recreate the entity.
\end{itemize}

Tombstones persist across compactions. They are \emph{reported} as eligible
under the watermark but \emph{retained} for delete safety.

\subsection{Notation field validation}

The server validates notation field keys against a whitelist defined in
\texttt{NotationMetadata}. Operations containing unknown notation keys are
rejected with \texttt{PRECONDITION\_FAILED}. This prevents clients from
emitting fields the server cannot merge deterministically.

The whitelist covers all currently supported notation fields for view objects
and connections. The notation inventory is kept in parity across:

\begin{itemize}
\item Server runtime whitelist constants (\texttt{ArchimeshService}).
\item JSON schema definitions (\texttt{schemas/ops.json}).
\item Client serializer/deserializer field coverage.
\item Client inventory test expectations \\
  (\texttt{OpMapperNotationInventoryTest}).
\end{itemize}

The script \texttt{scripts/check-notation-parity.sh} verifies this parity
and is wired into the preflight hardening checks.

\subsection{Lock model}

Lease-based locks are used for notation-sensitive edits:

\begin{itemize}
\item Operations requiring locks: \texttt{UpdateViewObjectOpaque},
  \texttt{UpdateConnectionOpaque}.
\item Lock targets are the affected view object or connection ids.
\item Default TTL: 10 seconds.
\item A client must hold the lock for a target before submitting an opaque
  notation update for that target.
\item If another client holds a conflicting lock, the submit is rejected
  with \texttt{LOCK\_CONFLICT}.
\item Locks are automatically released on WebSocket close.
\item Locks are advisory, not enforced by the merge layer --- they reduce
  destructive interleaving on noisy diagram edits while keeping semantic
  collaboration workable.
\end{itemize}

\subsection{CRDT maturity}

The CRDT implementation covers the currently supported operation set at the
P0 and P1 levels described in \filepath{crdt-spec.md}. Remaining work
includes:

\begin{itemize}
\item Closing the live relation-view sync gap where a semantic
  \texttt{CreateRelationship} can succeed but the corresponding
  \texttt{CreateConnection} is rejected or omitted.
\item Extending collection-aware operations to additional domains beyond
  property sets and view object children.
\item Gathering long-running workload evidence that metadata growth remains
  operationally bounded under production-like usage.
\end{itemize}

Any remaining non-CRDT domains are explicitly marked as unsupported rather
than left as implicit gaps.
